\documentclass[12pt,a4paper]{article}

% =========================
% Packages
% =========================
\usepackage[utf8]{inputenc}
\usepackage[T1]{fontenc}
\usepackage[margin=1in]{geometry}
\usepackage{times}
\usepackage{xcolor}
\usepackage{titlesec}
\usepackage{setspace}
\usepackage{enumitem}
\usepackage{graphicx}
\usepackage{float}
\usepackage{booktabs}
\usepackage{multirow}
\usepackage{array}
\usepackage{tabularx}
\usepackage{longtable}
\usepackage{amsmath}
\usepackage{amssymb}
\usepackage[hidelinks]{hyperref}
\usepackage{fancyhdr}

% =========================
% Formatting
% =========================
\definecolor{amritaMaroon}{HTML}{A4123F}
\titleformat{\section}{\Large\bfseries\color{amritaMaroon}}{\thesection}{1em}{}
\titleformat{\subsection}{\large\bfseries}{\thesubsection}{1em}{}
\titlespacing*{\section}{0pt}{10pt}{4pt}
\titlespacing*{\subsection}{0pt}{8pt}{3pt}
\onehalfspacing

\pagestyle{fancy}
\fancyhf{}
\fancyhead[L]{Project Phase 2 Report}
\fancyhead[R]{\thepage}
\renewcommand{\headrulewidth}{0.4pt}

\newcolumntype{L}{>{\raggedright\arraybackslash}X}
\newcolumntype{C}{>{\centering\arraybackslash}X}

\begin{document}

% =========================
% Cover Block
% =========================
\begin{center}
{\Large\textbf{\textcolor{amritaMaroon}{AMRITA VISHWA VIDYAPEETHAM}}}\\
{\large\textbf{School of Computing}}\\
{\large\textbf{Department of Computer Science and Engineering}}\\
\vspace{0.25cm}
{\LARGE\textbf{Project Phase 2 Report}}\\
\vspace{0.15cm}
{\large\textbf{Edge-Based AI System for Adaptive Regional Speech Enhancement in Autism Children}}\\
\vspace{0.15cm}
23CSE399 -- Panel Review 2\\
\today
\end{center}

\vspace{0.2cm}
\hrule
\vspace{0.3cm}

\begin{center}
\renewcommand{\arraystretch}{1.35}
\setlength{\tabcolsep}{5pt}
\begin{tabular}{|l|p{4.2cm}|p{4.2cm}|p{2.6cm}|c|}
\hline
\multicolumn{2}{|c|}{\textbf{Team Members}} & \textbf{Guide} & \textbf{Project Type} & \textbf{Marks} \\
\hline
CB.SC.U4CSE23018 & Rohith Kumar D &
\multirow{5}{=}{\raggedright \small \textbf{Dr. Venkataraman D} \\ \textit{Associate Professor, Dept. of CSE}} &
\multirow{5}{=}{\raggedright \small \textbf{Type 1} \\ Research} &
\multirow{5}{*}{\textbf{\Large \dots / 5}} \\ \cline{1-2}
CB.SC.U4CSE23055 & T Venkataramana & & & \\ \cline{1-2}
CB.SC.U4CSE23058 & VAB Jashwanth Reddy & & & \\ \cline{1-2}
CB.SC.U4CSE23060 & C Kalyan Kumar Reddy & & & \\ \cline{1-2}
CB.SC.U4CSE23446 & Hemanth Sholingaram & & & \\
\hline
\end{tabular}
\end{center}

\vspace{0.4cm}

\section*{Abstract}
This report presents the Phase 2 progression of a research project focused on building a privacy-first, edge-oriented speech therapy framework for children with Autism Spectrum Disorder (ASD), with strong emphasis on regional linguistic relevance for Tamil-speaking contexts. The core motivation is to provide clinically meaningful, objective, and scalable pronunciation support in environments where therapist access is limited and internet connectivity is unreliable.

Phase 1 established the conceptual architecture and broad research direction through an extensive literature survey spanning pronunciation scoring, adaptive therapy, multimodal analysis, and privacy-preserving AI. Panel feedback from Phase 1 identified feasibility risks in deploying a heavy, simultaneous multimodal stack on low-cost edge devices. In response, Phase 2 introduces a realistic staged design: local-first speech analysis, explicit edge-cloud boundaries, practical federated updates without expensive cryptographic overhead, and delayed RL activation until sufficient interaction data is collected.

The proposed technical pipeline combines Vocal Tract Length Normalization (VTLN), Wav2Vec 2.0 feature extraction, Dynamic Time Warping (DTW), and Goodness of Pronunciation (GOP) at syllable granularity, while tracking longitudinal articulatory trends through Vowel Space Area (VSA). A clear validation roadmap is defined with clinically relevant metrics, deployment constraints, and phased deliverables.

This document consolidates the Phase 1 findings, formalizes the cumulative gap, details the revised Phase 2 architecture, and provides a comprehensive execution plan toward pilot readiness and measurable therapeutic impact.

\newpage
\tableofcontents
\newpage

\section{Introduction}
\subsection{Background}
ASD is associated with communication and articulation challenges that demand repetitive, personalized, and timely intervention. In many regions of South India, access to trained speech therapists remains low relative to demand. This leads to interrupted therapy schedules, delayed progression, and heavy dependence on in-person specialist availability.

Existing digital therapy tools improve engagement but often remain clinically shallow in pronunciation diagnostics. Word-level correctness scores are commonly used, which are insufficient for therapist-level intervention because they do not localize which syllable or articulatory target failed. Further, many tools rely on cloud inference, creating privacy concerns for pediatric voice data and practical limitations in low-bandwidth settings.

\subsection{Need for Phase 2 Refinement}
The transition from proposal-level architecture to deployable research prototype requires rigorous technical prioritization. Phase 2 therefore focuses on narrowing the architecture to a feasible core that can run reliably on edge hardware while preserving therapeutic interpretability and data sovereignty.

\subsection{Scope of This Report}
This Phase 2 report includes:
\begin{itemize}[noitemsep,topsep=2pt]
\item Consolidated summary of findings from \textit{project\_report\_phase1.tex}.
\item Updated problem framing and originality statement.
\item Cumulative literature-driven gap formalization.
\item Detailed methodology and module-level pipeline.
\item Validation strategy, metrics, risks, and mitigation plan.
\item Current status, timeline, and next review deliverables.
\end{itemize}

\clearpage
\section{Phase 1 Findings Summary and Implications}
\subsection{Purpose of This Two-Page Summary}
This section provides a consolidated two-page summary of findings from the Phase 1 document and panel interaction. The objective is to preserve research continuity and show how Phase 2 design decisions were directly derived from validated observations rather than being introduced as disconnected changes.

Phase 1 produced valuable conceptual depth, but the panel correctly emphasized execution realism. The most important contribution of Phase 1 was therefore not only idea generation, but constraint discovery: what can be safely, efficiently, and clinically implemented first on low-cost edge infrastructure.

\subsection{What Phase 1 Established Successfully}
The Phase 1 report successfully established:
\begin{itemize}[noitemsep,topsep=2pt]
\item A strong societal and clinical motivation: therapist scarcity, intervention delay, and need for objective progression signals.
\item A broad literature-backed technical direction across speech biomarkers, adaptive learning, multimodal cues, and privacy-preserving architectures.
\item A clear intention to focus on Dravidian language challenges and atypical speech characteristics in ASD.
\item A clinically meaningful target of moving from coarse word-level evaluation to fine-grained, syllable-level insight.
\end{itemize}

These foundations remain valid in Phase 2 and are retained as core project principles.

\subsection{Critical Findings from Panel and Technical Review}
From the Phase 1 narrative and review feedback, the following findings were identified as high impact:
\begin{enumerate}[noitemsep,topsep=2pt]
\item \textbf{Overloaded Initial Stack:} The architecture attempted simultaneous deployment of heavy speech modeling, visual analysis, encryption-intensive privacy, and RL-based adaptation, which significantly increased system complexity and compute burden.
\item \textbf{Latency-Feasibility Conflict:} Claimed end-to-end latency targets were not compatible with the original full stack on intended edge hardware profiles.
\item \textbf{FHE Practicality Issue:} While cryptographically attractive, Full Homomorphic Encryption introduced unacceptable runtime overhead for real-time therapy flow.
\item \textbf{Premature RL Activation:} Day-1 adaptive RL was not clinically safe because policy quality depends on sufficient interaction episodes and stable reward shaping.
\item \textbf{Modeling Redundancy in Tracking Layer:} A heavy semantic graph layer was not required for early-phase progress tracking; simpler structured analytics were sufficient.
\item \textbf{Architecture Communication Gap:} The panel required clearer data-boundary representation between edge-resident processing and server-side aggregation.
\end{enumerate}

\subsection{Detailed Interpretation of Findings}
\textbf{Finding 1 and 2 (Complexity + Latency):} The combined effect of multi-module concurrency resulted in a high risk of unstable runtime behavior and poor user experience. In pediatric therapy settings, response delays reduce engagement and can distort behavioral signals. Therefore, Phase 2 prioritizes deterministic low-latency modules before adding optional advanced components.

\textbf{Finding 3 (FHE Overhead):} Privacy requirements are mandatory, but security design must remain deployable. The panel's concern was not about reducing privacy, but selecting an implementation path that preserves protection while avoiding impractical computational cost.

\textbf{Finding 4 (RL Timing):} Adaptive difficulty is desirable, but unsafe adaptivity can create inconsistent prompts and degrade therapeutic trust. Phase 2 therefore enforces a staged policy: rule-based control first, data accumulation second, RL onboarding third.

\textbf{Finding 5 and 6 (Tracking + Clarity):} For early clinical utility, clear reporting and understandable architecture are more valuable than structurally complex but under-validated data models. Phase 2 emphasizes clinician-readable outputs and explicit system boundaries.

\begin{table}[H]
\centering
\caption{Phase 1 Observation to Phase 2 Action Mapping}
\renewcommand{\arraystretch}{1.25}
\begin{tabularx}{\textwidth}{|p{3.7cm}|L|L|}
\hline
\textbf{Phase 1 Observation} & \textbf{Risk if Unchanged} & \textbf{Phase 2 Corrective Action} \\
\hline
Simultaneous heavy modules on edge & Runtime instability and delayed response & Phased module activation and priority ordering \\
\hline
Sub-200ms target under full stack & Unrealistic performance expectation & Core pipeline optimization before advanced add-ons \\
\hline
FHE-first privacy strategy & High compute overhead and scalability issues & Practical federated synchronization with strict local-first handling \\
\hline
Immediate RL adaptation & Unsafe policy behavior in low-data stage & Rule-based scaffold with RL readiness gating \\
\hline
Complex tracking model at startup & Increased engineering overhead, low near-term gain & Lightweight longitudinal tracking with structured metrics \\
\hline
Ambiguous architecture presentation & Review friction and implementation ambiguity & Explicit edge boundary and sync boundary documentation \\
\hline
\end{tabularx}
\end{table}

\newpage
\subsection{Net Impact of Phase 1 on Phase 2 Strategy}
Phase 1 findings produced a strategic shift from ``maximal architecture'' to ``validated incremental architecture.'' This shift improves scientific credibility because each subsystem now has a measurable purpose, a deployment rationale, and a validation criterion.

The Phase 2 engineering logic now follows four principles:
\begin{itemize}[noitemsep,topsep=2pt]
\item \textbf{Feasibility First:} Deploy what can be measured and stabilized on real target hardware.
\item \textbf{Clinical Actionability First:} Prioritize outputs that therapists can immediately interpret and use.
\item \textbf{Privacy by Deployable Design:} Preserve data sovereignty through practical local-first processing and controlled updates.
\item \textbf{Adaptive Intelligence by Readiness:} Enable RL only after data sufficiency and reward stability checks.
\end{itemize}

\subsection{Refined Phase Transition Statement}
Based on Phase 1 outcomes, the Phase 2 report adopts the following transition statement:

\textit{The project advances from broad conceptual ambition to a constrained, clinically interpretable, and hardware-realistic architecture. The Phase 2 system is intentionally staged to protect therapeutic reliability, ensure privacy compliance, and create robust baseline evidence before introducing high-variance adaptive intelligence components.}

\subsection{What This Means for Evaluation in Phase 2}
The revised strategy also changes how success is measured. Instead of evaluating many loosely integrated modules at once, Phase 2 evaluates tightly-coupled core objectives:
\begin{enumerate}[noitemsep,topsep=2pt]
\item Can syllable-level outputs remain stable and interpretable across repeated sessions?
\item Can the pipeline maintain acceptable runtime on intended edge constraints?
\item Can privacy-preserving data handling be sustained without workflow disruption?
\item Can progression logic remain safe and transparent before RL activation?
\end{enumerate}

\subsection{Summary of Phase 1-to-Phase 2 Learning}
In summary, Phase 1 served as a discovery stage that surfaced execution bottlenecks early. The major value of that stage was not merely novelty exploration, but informed reduction of technical risk. The resulting Phase 2 plan is therefore stronger: it retains the original clinical goal, keeps research originality at syllable granularity and regional adaptation, and introduces an implementation path that is realistic for rural edge deployment scenarios.

\clearpage

\section{Problem Definition and Originality}
\subsection{Refined Problem Statement}
The project aims to build an offline-capable, privacy-first, edge-based speech therapy system for autistic children that can evaluate pronunciation at syllable level, adapt to regional phonetic realities, and support objective longitudinal tracking for therapists and caregivers.

\subsection{Originality Claim}
This work is an integrated extension over existing literature with originality in composition and deployment strategy. The novelty is not in a single isolated algorithm, but in the clinically grounded combination of:
\begin{itemize}[noitemsep,topsep=2pt]
\item VTLN for pediatric acoustic normalization,
\item DTW+GOP for syllable-precise scoring,
\item staged adaptive difficulty control,
\item practical federated model evolution,
\item regional phonetic orientation for Dravidian articulation patterns.
\end{itemize}

\subsection{Research Questions}
\begin{itemize}[noitemsep,topsep=2pt]
\item Can syllable-level scoring provide clearer therapy guidance than word-level scoring?
\item Can a local-first architecture preserve privacy without sacrificing usability?
\item Can staged adaptivity outperform static progression under rural deployment constraints?
\item Can proxy disordered-speech resources bridge data scarcity until primary corpus collection scales?
\end{itemize}

\section{Literature Survey and Cumulative Gap Analysis}
\subsection{Survey Coverage}
The Phase 1 survey included 25 papers spanning pronunciation scoring, emotion and multimodal channels, digital therapy interfaces, transfer learning for low-resource languages, adaptive RL, and federated privacy architectures.

\subsection{Thematic Synthesis}
\textbf{Theme A: Clinical Precision}\
Most systems score speech at utterance or word level. Few provide syllable-level error localization tied to intervention.

\textbf{Theme B: Regional-Linguistic Adaptation}\
High-resource language dominance leads to weak performance for Dravidian phonetic characteristics and atypical pediatric prosody.

\textbf{Theme C: Privacy and Deployment}\
Cloud dependency increases data-exposure risk and limits continuity in low-connectivity locations.

\textbf{Theme D: Adaptivity}\
RL offers promise, but safe deployment requires phased onboarding and robust reward design.

\subsection{Cumulative Gap Statement}
There is no integrated rural-ready system that simultaneously satisfies:
\begin{itemize}[noitemsep,topsep=2pt]
\item strict local-first processing,
\item syllable-level pronunciation analysis,
\item pediatric acoustic normalization,
\item staged adaptive therapy logic,
\item objective longitudinal metrics usable by clinicians.
\end{itemize}

\begin{table}[H]
\centering
\caption{Gap Mapping from Phase 1 Survey to Phase 2 Design}
\renewcommand{\arraystretch}{1.25}
\begin{tabularx}{\textwidth}{|p{3.2cm}|L|L|}
\hline
\textbf{Gap Area} & \textbf{Observed Limitation in Existing Systems} & \textbf{Phase 2 Response} \\
\hline
Precision & Word-level or coarse scoring; low therapeutic actionability & DTW + GOP at syllable level \\
\hline
Regionality & Weak adaptation for Dravidian articulation patterns & Two-stage transfer strategy + regional targets \\
\hline
Privacy & Cloud-heavy inference pipelines & Edge-first processing and controlled sync \\
\hline
Adaptivity & Static lesson progression in many tools & Rule-based now, DDPG when data-safe \\
\hline
Progress tracking & Subjective or sparse reporting & VSA/GOP longitudinal trajectory tracking \\
\hline
\end{tabularx}
\end{table}

\section{System Objectives and Success Criteria}
\subsection{Primary Objectives}
\begin{enumerate}[noitemsep,topsep=2pt]
\item Build edge-first pronunciation evaluation pipeline.
\item Ensure clinically interpretable syllable-level outcomes.
\item Preserve privacy through local-first data handling.
\item Establish reliable therapist-facing progress signals.
\item Prepare adaptive control transition under safe conditions.
\end{enumerate}

\subsection{Success Criteria}
\begin{itemize}[noitemsep,topsep=2pt]
\item Demonstrable reduction in baseline recognition error.
\item Positive clinician-aligned reliability for GOP scoring.
\item Consistent latency under target device constraints.
\item Stable generation of longitudinal progress summaries.
\item Documented transition plan from rule-based to RL.
\end{itemize}

\section{Proposed Methodology}
\subsection{Pipeline Overview}
The proposed architecture is a layered dataflow:
\begin{enumerate}[noitemsep,topsep=2pt]
\item \textbf{Capture}: Pediatric audio input with session metadata.
\item \textbf{Normalize}: VTLN to account for vocal tract differences.
\item \textbf{Encode}: Wav2Vec-derived representation for robust feature extraction.
\item \textbf{Align}: DTW for variable-rate temporal alignment.
\item \textbf{Score}: GOP and related confidence measures.
\item \textbf{Track}: VSA and historical performance trends.
\item \textbf{Adapt}: Rule-based progression; RL gated for later stage.
\item \textbf{Sync}: Aggregated updates and analytics under controlled policies.
\end{enumerate}

\subsection{Core Algorithms}
\textbf{VTLN:} compensates for child-adult acoustic mismatch.\
\textbf{DTW:} aligns sequences with speaking-rate variation.\
\textbf{GOP:} evaluates pronunciation confidence at sub-word units.\
\textbf{VSA:} supports longitudinal articulatory-motor interpretation.

\subsection{Conceptual Flow Figure}
\begin{figure}[H]
\centering
\fbox{\begin{minipage}[c][5cm]{0.9\textwidth}
\centering
\vspace{1cm}
\textbf{Phase 2 Processing Flow}\\
\vspace{0.2cm}
Capture $\rightarrow$ VTLN $\rightarrow$ Representation $\rightarrow$ DTW\\
$\rightarrow$ GOP/VSA $\rightarrow$ Difficulty Controller $\rightarrow$ Tracking/Reporting
\end{minipage}}
\caption{Conceptual methodology used in Phase 2}
\end{figure}

\section{Module-Wise Design}
\subsection{Module 1: Data Acquisition and Session Control}
Handles utterance recording, lesson context association, and repeat-attempt management.

\subsection{Module 2: Acoustic Normalization and Representation}
Performs normalization and feature extraction robust to speaker variability and atypical speech rhythms.

\subsection{Module 3: Alignment and Pronunciation Scoring}
Produces syllable-level diagnostics suitable for clinician interpretation.

\subsection{Module 4: Progress Tracking and Reporting}
Maintains longitudinal trend records and prepares structured reports for therapist review.

\subsection{Module 5: Adaptive Difficulty Logic}
Executes transparent rule-based adaptation during data-scarce stage and defines RL transition guardrails.

\subsection{Module 6: Privacy and Sync Boundary}
Enforces local-first processing and controlled synchronization events for model and analytics evolution.

\section{Datasets and Data Strategy}
\subsection{Constraint}
Primary Tamil ASD disordered speech data availability is limited, requiring staged data strategy.

\subsection{Proxy + Regional Strategy}
Phase 2 uses a combination of:
\begin{itemize}[noitemsep,topsep=2pt]
\item regional language corpora for phonetic grounding,
\item disordered speech proxy datasets for anomaly adaptation,
\item planned primary corpus collection under ethical controls.
\end{itemize}

\begin{table}[H]
\centering
\caption{Dataset Utilization Strategy}
\renewcommand{\arraystretch}{1.25}
\begin{tabularx}{\textwidth}{|p{3.5cm}|L|L|}
\hline
\textbf{Dataset Category} & \textbf{Role in Pipeline} & \textbf{Phase 2 Usage} \\
\hline
Regional speech corpora & Phonetic adaptation and baseline modeling & Stage-1 adaptation \\
\hline
Disordered speech proxy & Robustness to atypical articulation patterns & Stage-2 adaptation \\
\hline
Primary ASD corpus (planned) & Clinical validation with target population & Pilot preparation \\
\hline
\end{tabularx}
\end{table}

\section{Validation and Evaluation Plan}
\subsection{Quantitative Metrics}
\begin{itemize}[noitemsep,topsep=2pt]
\item Recognition error reduction relative to baseline.
\item GOP-clinician agreement trend.
\item Latency and throughput under edge constraints.
\item Consistency of syllable-boundary scoring.
\item Longitudinal improvement indications in VSA trajectory.
\end{itemize}

\subsection{Qualitative Metrics}
\begin{itemize}[noitemsep,topsep=2pt]
\item Therapist interpretability of per-session output.
\item Caregiver usability and session continuity feedback.
\item Child engagement stability over repeated sessions.
\end{itemize}

\begin{table}[H]
\centering
\caption{Validation Matrix}
\renewcommand{\arraystretch}{1.25}
\begin{tabularx}{\textwidth}{|p{3.5cm}|L|L|}
\hline
\textbf{Metric Group} & \textbf{Target Direction} & \textbf{Validation Method} \\
\hline
Accuracy & Improve over baseline & Controlled benchmark comparison \\
\hline
Interpretability & Increase therapist usefulness & Expert review and rubric scoring \\
\hline
Latency & Within edge limits & Device profiling under standard workload \\
\hline
Adaptivity safety & Avoid unstable behavior & Rule-gated progression checks \\
\hline
\end{tabularx}
\end{table}

\section{Risk Analysis and Mitigation}
\subsection{Technical Risks}
\begin{itemize}[noitemsep,topsep=2pt]
\item Model drift due to heterogeneous speaking patterns.
\item Overfitting to proxy datasets.
\item Edge performance bottlenecks in low-cost hardware.
\end{itemize}

\subsection{Clinical and Operational Risks}
\begin{itemize}[noitemsep,topsep=2pt]
\item Misinterpretation of automated scores without contextual review.
\item Inconsistent session conditions affecting score comparability.
\item Premature RL deployment causing unstable task adaptation.
\end{itemize}

\subsection{Mitigation Strategy}
\begin{itemize}[noitemsep,topsep=2pt]
\item Staged rollout with explicit readiness gates.
\item Hybrid review process with therapist-in-the-loop checkpoints.
\item Standardized recording guidance and quality controls.
\item Incremental optimization and device-specific profiling.
\end{itemize}

\section{Ethical, Privacy, and Governance Considerations}
\subsection{Privacy Principles}
\begin{itemize}[noitemsep,topsep=2pt]
\item Local-first processing of sensitive pediatric voice data.
\item Minimal required synchronization and data minimization.
\item Explicit separation of analytics payloads from raw sensitive artifacts.
\end{itemize}

\subsection{Ethical Use Principles}
\begin{itemize}[noitemsep,topsep=2pt]
\item The system assists therapy; it does not replace clinical judgment.
\item Outputs must remain interpretable and contestable.
\item Pilot usage should align with institutional ethical review processes.
\end{itemize}

\section{Implementation Timeline and Work Breakdown}
\subsection{Phase-Wise Timeline}
\begin{itemize}[noitemsep,topsep=2pt]
\item \textbf{Window A:} Core scoring pipeline stabilization.
\item \textbf{Window B:} Optimization and reporting integration.
\item \textbf{Window C:} Pilot preparation and reliability assessment.
\item \textbf{Window D:} Adaptive upgrade readiness evaluation.
\end{itemize}

\begin{table}[H]
\centering
\caption{Work Breakdown for Next Cycle}
\renewcommand{\arraystretch}{1.25}
\begin{tabularx}{\textwidth}{|p{4cm}|L|p{2.8cm}|}
\hline
\textbf{Work Package} & \textbf{Output} & \textbf{Status Focus} \\
\hline
Syllable scoring pipeline & Stable DTW+GOP flow & In progress \\
\hline
Longitudinal tracker & Session-wise progress summary & In progress \\
\hline
Edge optimization & Measured runtime profile & Planned \\
\hline
Pilot protocol package & Validation checklist and forms & Planned \\
\hline
Adaptive upgrade gating & RL readiness criteria document & Planned \\
\hline
\end{tabularx}
\end{table}

\section{Current Status and Next Deliverables}
\subsection{Current Status}
\begin{itemize}[noitemsep,topsep=2pt]
\item Problem framing and architecture constraints are now explicit.
\item Gap analysis is complete and aligned to module design.
\item Methodology and validation matrix are defined for execution.
\item Risk and ethics sections are formalized for pilot readiness.
\end{itemize}

\subsection{Deliverables for Next Review}
\begin{itemize}[noitemsep,topsep=2pt]
\item End-to-end demonstrable scoring workflow.
\item Initial benchmark results and latency profile.
\item Therapist-readable progress output samples.
\item Refined adaptive strategy with transition criteria.
\end{itemize}

\section{Conclusion}
Phase 2 transforms the project from concept-heavy architecture to an execution-ready, clinically interpretable, and deployment-aware research pipeline. The revised design respects hardware constraints, privacy obligations, and therapy realities while maintaining scientific ambition through syllable-level diagnostics and adaptive planning. The resulting framework is positioned for structured pilot validation and measurable impact in underserved settings.

\section*{References (Consolidated from Phase 1/Phase 2 Sources)}
\begin{enumerate}[leftmargin=6mm,itemsep=2pt,topsep=2pt]
\item J. Zhu, Y. Li, and X. Zhang, ``The taste of IPA: Towards open-vocabulary keyword spotting,'' IEEE ICASSP, 2024.
\item A. Landowska et al., ``Automatic Emotion Recognition in Children with Autism: A Systematic Review,'' Sensors, 2022.
\item S. Saeedi et al., ``Application of Digital Games for Speech and Language Therapy in Children,'' Journal of Biomedical Informatics, 2022.
\item M. S. Remya et al., ``Artificial Intelligence Techniques for Speech Disorder Classification: A Review,'' IEEE Access, 2025.
\item S. Sandoval et al., ``Automatic assessment of vowel space area,'' Journal of the Acoustical Society of America, 2013.
\item B. Bharathi et al., ``Speech recognition for vulnerable individuals in Tamil,'' LT-EDI Workshop, 2023.
\item A. Girirajan et al., ``Real-time Tamil ASR using Wav2Vec2,'' 2024.
\item T. Sharma et al., ``Edge AI: technologies, applications, and challenges,'' ACET, 2024.
\item C.-H. Chueh et al., ``Generative AI-powered clinical language therapy,'' JAMA Pediatrics, 2025.
\item N. Pavlidis et al., ``Federated learning for privacy-preserving anomaly detection in ASD diagnosis,'' IEEE TIFS, 2024.
\end{enumerate}

\end{document}
